\documentclass{article}
\usepackage{amsmath,amssymb}
\usepackage[vocab={orthogonalpolynomials}]{leanify}
\newtheorem{lemma}{Lemma}
\begin{document}

\section{Notation}

We write $P_n$ for the Legendre polynomial of degree $n$.
We write $H_n$ for the Hermite polynomial of degree $n$.
We write $L_n$ for the Laguerre polynomial of degree $n$.
We write $T_n$ for the Chebyshev polynomial of the first kind of degree $n$.
We write $U_n$ for the Chebyshev polynomial of the second kind of degree $n$.

\section{Values}

\begin{lemma}\label{legendre}
$P_2(x)=\frac{3x^2-1}{2}$ for every real number $x$, and $P_n(1)=1$ for every nonnegative integer $n$.
\end{lemma}

\begin{lemma}\label{hermite}
For every real number $x$, $H_2(x)=4x^2-2$ and $L_1(x)=1-x$.
\end{lemma}

\begin{lemma}\label{chebyshev}
For every real number $x$, $T_2(x)=2x^2-1$ and $U_2(x)=4x^2-1$. For every nonnegative integer $n$, $T_n(1)=1$.
\end{lemma}

\section{Explicit sums}

\begin{lemma}\label{sums}
For every nonnegative integer $n$ and every real number $x$, $T_n(x)=\sum_{k=0}^{\lfloor n/2\rfloor}\binom{n}{2k}(x^2-1)^kx^{n-2k}$ and $P_n(x)=\sum_{k=0}^{n}\binom{n}{k}^2(\frac{x-1}{2})^{n-k}(\frac{x+1}{2})^k$.
\end{lemma}

\end{document}
